\section{Solving one-dimensional potentials (10/2)}

\referto{\tong~2.1.1, 2.3.2, 2.4.1; \sakurai~Appendix B on ``Elementary Solutions to Schr\"odinger Wave Equation"}

\subsection{General Properties of stationary Schr\"odinger equations}

The stationary Schr\"odinger equation is
\begin{equation}
 \dbra{x}\hat H\dket{\psi}=E\expval{x}{\psi}.
\end{equation}
For a particle of mass $m$ in a one-dimensional potential $V(x)$,
\begin{equation}
 \boxed{
 -\frac{\hbar^2}{2m}\psi''(x)+V(x)\psi(x)=E\psi(x).}
\end{equation}
This is a linear, second-order ordinary differential equation.

\textbf{Linearity:}

If $\psi_1(x)$ and $\psi_2(x)$ are solutions with the same energy $E$, then
any linear combination
\begin{equation}
 c_1\psi_1(x)+c_2\psi_2(x),
 \qquad c_1,c_2\in\mathbb C,
\end{equation}
is also a solution with energy $E$.

\textbf{Complex conjugation and time reversal symmetry}

If $V(x)$ is real and $\psi(x)$ is a solution with energy $E$, then
$\psi^*(x)$ is also a solution with the same energy. For
\begin{equation}
 \hat H=\frac{\hat p^{\,2}}{2m}+V(\hat x),
\end{equation}
this follows from
\begin{equation}
 \hat H^*=\hat H.
\end{equation}
This is the stationary manifestation of time-reversal symmetry for a
spinless particle in a real scalar potential.

\textbf{Real eigenfunctions}

Energy eigenfunctions can be chosen to be real. Starting from a complex
solution $\psi(x)$, define
\begin{equation}
 \psi_1(x)=\psi(x)+\psi^*(x),
 \qquad
 \psi_2(x)=\imth\bigl(\psi(x)-\psi^*(x)\bigr).
\end{equation}
Both $\psi_1$ and $\psi_2$ are real solutions with the same energy. A
nonzero member of this pair can therefore be used as a real eigenfunction.

\textbf{Smoothness and matching conditions}

If $V(x)$ is smooth, then solutions of the Schr\"odinger equation are smooth.
If $V(x)$ has a finite discontinuity, integration of the equation across the
discontinuity shows that both $\psi(x)$ and $\psi'(x)$
are continuous. Singular potentials, such as a delta function, require a
modified derivative matching condition.

\subsection{Examples}

\subsubsection{Free particle}
For a free particle,
\begin{equation}
 V(x)=0.
\end{equation}
The plane-wave solutions are
\begin{equation}
 \psi_k(x)\propto e^{\imth kx},
 \qquad
 E=\frac{\hbar^2k^2}{2m}.
\end{equation}
These are scattering states and are not square-normalizable on the entire
real line. The two waves $\psi_{+k}$ and $\psi_{-k}$ have the same energy.
Real eigenfunctions may be chosen as $\cos(kx)$ and $\sin(kx)$. 

To properly normalize these plane waves, we have to utilize the Dirac delta function:
\begin{align}
    \expval{k_1|k_2}=\int\dif{x}~\psi_{k_1}^\ast(x)\psi_{k_2}(x)=\delta(k_1-k_2)
\end{align}
which means
\begin{align}
    \psi_k(x)=\frac1{\sqrt{2\pi}}e^{\imth kx}
\end{align}


\textbf{Periodic boundary condition:}

If we consider a plane wave on a periodic ring of length $L$, we need to enforce the \emph{periodic boundary condition} 
\begin{align}
    \psi_k(x)\equiv\psi_k(x+L)
\end{align}
which quantizes the wavevector $k$ as 
\begin{align}
    k=\frac{2\pi}{L}n_k,~~~n_k\in\mbz.
\end{align}


\subsubsection{Step potential}
Consider
\begin{equation}
 V(x)=
 \begin{cases}
 0, & x<0,\\
 U, & x>0,
 \end{cases}
 \qquad U>0.
\end{equation}
Because $V(x)\geq0$, there are no energy eigenstates with $E<0$.

\textbf{Case I: $E>U$}

Define
\begin{equation}
 k=\frac{\sqrt{2mE}}{\hbar},
 \qquad
 k'=\frac{\sqrt{2m(E-U)}}{\hbar}.
\end{equation}
For a wave incident from the left, write
\begin{equation}
 \psi(x)=
 \begin{cases}
 e^{\imth kx}+r e^{-\imth kx}, & x<0,\\
 t e^{\imth k'x}, & x>0.
 \end{cases}
\end{equation}
Continuity of $\psi$ and $\psi'$ at $x=0$ gives
\begin{equation}
 1+r=t,
 \qquad
 k(1-r)=k't.
\end{equation}
Therefore,
\begin{equation}
 \boxed{r=\frac{k-k'}{k+k'},
 \qquad
 t=\frac{2k}{k+k'}=\frac{2}{1+k'/k}.}
\end{equation}
The reflection and transmission probabilities are determined by probability
currents:
\begin{equation}
 R=|r|^2,
 \qquad
 T=\frac{k'}{k}|t|^2,
 \qquad
 R+T=1.
\end{equation}

\textbf{Case II: $0<E<U$}

Now define
\begin{equation}
 k=\frac{\sqrt{2mE}}{\hbar},
 \qquad
 \kappa=\frac{\sqrt{2m(U-E)}}{\hbar},
 \qquad k'=\imth\kappa.
\end{equation}
The physical solution is
\begin{equation}
 \psi(x)=
 \begin{cases}
 e^{\imth kx}+r e^{-\imth kx}, & x<0,\\
 t e^{-\kappa x}, & x>0.
 \end{cases}
\end{equation}
The growing exponential for $x>0$ is excluded. The matching formulas are the
same as before after substituting $k'=\imth\kappa$:
\begin{equation}
 t=\frac{2k}{k+\imth\kappa},
 \qquad
 r=\frac{k-\imth\kappa}{k+\imth\kappa}.
\end{equation}
Thus $|r|^2=1$. There is no transmitted probability current, although the
wave function penetrates a distance of order $\kappa^{-1}$ into the
classically forbidden region.

\subsubsection{Delta-function potential}
Consider 
\begin{equation}
 V(x)=D_0~\delta(x).
\end{equation}
The wave function is continuous at the origin. Integrating the stationary
Schr\"odinger equation from $-\epsilon$ to $+\epsilon$ and taking
$\epsilon\to0^+$ gives
\begin{equation}
 \boxed{
 \psi'(0^+)-\psi'(0^-)=\frac{2mD_0}{\hbar^2}\psi(0).}
\end{equation}

\textbf{Scattering states: $E>0$}

Let
\begin{equation}
 k=\frac{\sqrt{2mE}}{\hbar},
\end{equation}
and use
\begin{equation}
 \psi(x)=
 \begin{cases}
 e^{\imth kx}+r e^{-\imth kx}, & x<0,\\
 t e^{\imth kx}, & x>0.
 \end{cases}
\end{equation}
Continuity gives
\begin{equation}
 1+r=t.
\end{equation}
The derivative jump condition gives
\begin{equation}
 \imth k(t-1+r)=\frac{2mD_0}{\hbar^2}t.
\end{equation}
Define
\begin{equation}
 \alpha=\frac{mD_0}{\hbar^2k}.
\end{equation}
Then
\begin{equation}
 \boxed{t=\frac{1}{1+\imth\alpha},
 \qquad
 r=\frac{-\imth\alpha}{1+\imth\alpha}.}
\end{equation}
Equivalently,
\begin{equation}
 t=\frac{1}{1-mD_0/(\imth k\hbar^2)},
 \qquad
 r=\frac{mD_0/(\imth k\hbar^2)}{1-mD_0/(\imth k\hbar^2)}.
\end{equation}
The fluxes satisfy
\begin{equation}
 |t|^2+|r|^2=1.
\end{equation}
The same scattering result applies to either sign of $D_0$ as long as
$E>0$.

\textbf{Bound state for an attractive delta potential}

Let $D_0<0$ and $E<0$. Define
\begin{equation}
 \kappa=\frac{\sqrt{2m|E|}}{\hbar}.
\end{equation}
Normalizability requires exponential decay at both infinities:
\begin{equation}
 \psi(x)=
 \begin{cases}
A e^{\kappa x}, & x<0,\\
 B e^{-\kappa x}, & x>0.
 \end{cases}
\end{equation}
The continuity gives $A=B$ and the derivative jump condition becomes
\begin{equation}
 -2\kappa A=\frac{2mD_0}{\hbar^2}A,
\end{equation}
so
\begin{equation}
 \kappa=-\frac{mD_0}{\hbar^2}=\frac{m|D_0|}{\hbar^2}.
\end{equation}
Therefore the attractive delta potential has one bound state with energy
\begin{equation}
 \boxed{E=-\frac{mD_0^2}{2\hbar^2}.}
\end{equation}
Normalization gives
\begin{equation}
 1=\int_{-\infty}^{\infty}|\psi(x)|^2\dif x
 =\frac{|A|^2}{\kappa},
\end{equation}
so one may choose
\begin{equation}
 \boxed{\psi(x)=\sqrt{\kappa}\,e^{-\kappa|x|}.}
\end{equation}

\subsection{Bound and Scattering Spectra}

Bound states are normalizable and have discrete energies. For the attractive
delta potential, there is a single bound-state energy. In contrast, scattering states are
not square-normalizable in the full space and form a continuous energy
spectrum. They are instead normalized using delta-function normalization. 



\subsection{Summary}
%\addcontentsline{toc}{chapter}{Summary}
\begin{itemize}
 \item The stationary Schr\"odinger equation is a linear, second-order ODE.
 \item For real potentials, energy eigenfunctions can be chosen real.
 \item Finite potential discontinuities preserve continuity of both $\psi$
       and $\psi'$.
 \item A delta potential preserves $\psi$ but produces a finite jump in
       $\psi'$.
 \item Step and delta potentials illustrate reflection, transmission,
       evanescent waves, and bound-state quantization.
 \item Bound states have discrete energies, whereas scattering states form a
       continuous spectrum.
\end{itemize}

